\section{Conclusion}\label{sec:conclusion}

This paper introduced DSPy, a new programming model for designing AI systems using pipelines of pretrained LMs and other tools. We presented three new concepts introduced in this abstraction (DSPy signatures, modules, and teleprompters), and showed in two very different case studies that it supports rapid development of highly effective systems that use relatively small LMs. We have maintained open-source versions of this framework for close to a year. In this period, we have seen and created a large number of programs that were compiled to high-quality systems by DSPy, spanning tasks from information extraction to low-resource synthetic data generation. In the interest of space and to maintain reasonable scope in this paper, we leave reporting on such tasks under controlled experimental conditions to future work. While in-context learning has proved transformative over the past 2--3 years of LM research, we argue that the true expressive power in this emerging paradigm is in building sophisticated text transformation graphs in which composable modules and optimizers (teleprompters) come together to leverage LMs in more systematic and reliable ways.


